\section{Поточний стан досліджень}

Задачі аналізу музичних даних активно досліджуються у сфері інтелектуального аналізу даних та машинного навчання. Одним із ключових напрямків є дослідження аудіофіч музичних композицій для побудови рекомендаційних систем, прогнозування популярності треків та автоматичної класифікації жанрів.

У сучасних музичних сервісах використовуються алгоритми машинного навчання для персоналізації рекомендацій та аналізу поведінки користувачів. Spotify застосовує аудіофічі, які формуються на основі цифрової обробки сигналів та статистичного аналізу музичних характеристик.

У багатьох дослідженнях використовуються методи кластеризації для групування композицій за схожими ознаками. Найбільш поширеним алгоритмом є KMeans, який дозволяє ефективно виявляти приховані структури у великих наборах даних.

Також активно застосовуються методи класифікації для прогнозування популярності треків. Для цього використовуються алгоритми Logistic Regression, Random Forest, Gradient Boosting та нейронні мережі. Практика показує, що ансамблеві методи часто демонструють кращі результати для музичних даних через нелінійні взаємозв'язки між ознаками.

У даній роботі використовується комбінація класичних методів аналізу даних та машинного навчання, що дозволяє побудувати інтерпретоване та відтворюване дослідження без використання складних глибоких нейронних мереж.
